% RQ4 Theoretical Framework: Treasury Resilience

We develop formal models of treasury dynamics and resilience metrics.

\subsection{Treasury Dynamics Model}

\subsubsection{Treasury State}

Treasury state at time $t$:

\begin{equation}
T(t) = T(t-1) + I(t) - O(t) + \Delta V(t)
\end{equation}

where:
\begin{itemize}
    \item $I(t)$ = inflows (revenue, contributions)
    \item $O(t)$ = outflows (proposal spending, operations)
    \item $\Delta V(t)$ = value change (market movements affecting token holdings)
\end{itemize}

\subsubsection{Buffer Mechanics}

With buffer fraction $b$:

\begin{equation}
\text{available}(t) = (1 - b) \cdot T(t)
\end{equation}

\begin{equation}
\text{buffer}(t) = b \cdot T(t)
\end{equation}

Spending is constrained: $O(t) \leq \text{available}(t)$

\subsubsection{Emergency Top-Up}

When buffer falls below threshold $\tau_{\text{buffer}}$:

\begin{equation}
\text{top-up}(t) = \mathbf{1}[\text{buffer}(t) < \tau_{\text{buffer}}] \cdot \min(\text{deficit}, \text{top-up rate})
\end{equation}

where $\text{deficit} = \tau_{\text{buffer}} - \text{buffer}(t)$

Top-up sources reduce available spending or divert inflows.

\subsection{Resilience Metrics}

\subsubsection{Treasury Volatility}

Standard deviation of treasury value over rolling window:

\begin{equation}
\sigma_T = \sqrt{\frac{1}{w} \sum_{i=t-w}^{t} (T(i) - \bar{T})^2}
\end{equation}

Lower volatility indicates more stable treasury.

\subsubsection{Drawdown Risk}

Maximum peak-to-trough decline:

\begin{equation}
\text{MaxDrawdown} = \max_{t} \left( \frac{\max_{s \leq t} T(s) - T(t)}{\max_{s \leq t} T(s)} \right)
\end{equation}

\subsubsection{Runway}

Time until treasury depletion at current burn rate:

\begin{equation}
\text{Runway} = \frac{T(t)}{\bar{O}}
\end{equation}

where $\bar{O}$ is average outflow rate.

\subsubsection{Buffer Coverage}

Ratio of buffer to expected outflows:

\begin{equation}
\text{Coverage} = \frac{\text{buffer}(t)}{\mathbb{E}[O]}
\end{equation}

Higher coverage indicates better protection against outflow shocks.

\subsection{Tradeoff Analysis}

\subsubsection{Resilience vs. Efficiency}

Buffer reserves create tradeoff:
\begin{itemize}
    \item \textbf{High buffer}: More resilient but idle capital
    \item \textbf{Low buffer}: Higher efficiency but vulnerable to shocks
\end{itemize}

Optimal buffer depends on:
\begin{itemize}
    \item Inflow volatility
    \item Outflow predictability
    \item Market conditions (treasury value volatility)
    \item Risk tolerance
\end{itemize}

\subsubsection{Spending Limits vs. Operational Flexibility}

Spending limits protect treasury but may:
\begin{itemize}
    \item Block valuable proposals exceeding limits
    \item Fragment large initiatives into multiple smaller proposals
    \item Create waiting queues for spending approval
\end{itemize}

\subsection{Theoretical Predictions}

\subsubsection{Prediction 1: Buffer-Volatility Relationship}

Treasury volatility decreases with buffer fraction:

\begin{equation}
\sigma_T \propto (1 - b)
\end{equation}

as buffer dampens spending-driven volatility.

\subsubsection{Prediction 2: Spending Limit-Runway Relationship}

Tighter spending limits extend runway:

\begin{equation}
\text{Runway} \propto \frac{1}{\text{spending\_limit}}
\end{equation}

\subsubsection{Prediction 3: Top-Up Effectiveness}

Emergency top-ups reduce drawdown severity when:
\begin{itemize}
    \item Inflow capacity exceeds top-up requirements
    \item Top-up triggers before severe depletion
    \item Recovery rate exceeds continued outflow rate
\end{itemize}

\subsection{Market Condition Modeling}

We model treasury value volatility as:

\begin{equation}
\Delta V(t) = T_{\text{token}}(t) \cdot r(t)
\end{equation}

where $r(t) \sim \mathcal{N}(\mu, \sigma_{\text{market}}^2)$ is the token return.

Different market regimes:
\begin{itemize}
    \item \textbf{Bull}: $\mu > 0$, moderate $\sigma$
    \item \textbf{Bear}: $\mu < 0$, high $\sigma$
    \item \textbf{Stable}: $\mu \approx 0$, low $\sigma$
\end{itemize}
